% ECONOMICS_PROBLEM: {"id": "I38-454", "title": "Scaling durable graduation programs", "jel_code": "I38", "source_id": "OP-0371"}
% !TeX program = xelatex
\documentclass[11pt,a4paper]{article}
\usepackage[margin=23mm]{geometry}
\usepackage{fontspec}
\setmainfont{Latin Modern Roman}
\usepackage{amsmath,amssymb,mathtools}
\usepackage{xcolor,enumitem,booktabs,longtable,array,xurl,hyperref}
\definecolor{muted}{HTML}{53636C}
\hypersetup{colorlinks=true,urlcolor=blue,linkcolor=blue}
\setlength{\parindent}{0pt}
\setlength{\parskip}{5pt}
\newcommand{\R}{\mathbb R}
\newcommand{\E}{\mathbb E}
\newcommand{\N}{\mathbb N}
\newcommand{\ind}{\mathbf 1}
\newcommand{\Var}{\operatorname{Var}}
\newcommand{\Cov}{\operatorname{Cov}}
\newcommand{\supp}{\operatorname{supp}}
\DeclareMathOperator*{\argmin}{arg\,min}
\DeclareMathOperator*{\argmax}{arg\,max}
\newcommand{\entrymeta}[3]{{\sffamily\small\color{muted}\textbf{\texttt{#1}}\quad|\quad #2\par #3\par}}
\newcommand{\Problemref}[1]{\texttt{#1}}
\newcommand{\JELref}[2]{\texttt{#2} (JEL \texttt{#1})}
\begin{document}
\section*{Scaling durable graduation programs}
\textbf{Problem ID:} \texttt{I38-454}\quad \textbf{JEL:} \texttt{I38}\par
% BEGIN ECONOMICS STATEMENT
\label{problem:OP-0371}
{\sffamily\small\color{muted}Original subject group: Development and poverty\par}
\label{definitions:problem:30007653}
\textbf{Audit classification:} Published research agenda. Primary institutional agenda checked. Component doses, delivery and population adaptations are explicit questions; formal objectives and noninferiority criteria are editorial. Page is a living resource.
\subsection*{Definitions and precise research target}
\textbf{Complete editorial problem.} Fix a target population and a government delivery system for an economic-inclusion program. Let \(p=(a,t,c,s)\) specify the quantities of productive assets, training, coaching, and savings support. Let \(q\in[0,1]\) be coverage, allowing implementation quality, costs, market prices, and spillovers to change with scale. Household \(i\)'s potential consumption at horizon \(T\) is \(C_{iT}(p,q)\). Define \(\tau_T(p,q)=E[C_{iT}(p,q)-C_{iT}(0,0)]\) for the entire target population, separately reporting beneficiaries and nonbeneficiaries. Let \(K(p,q)\) be total fiscal and implementation cost, and let \(B\) be an explicit budget. Characterize combinations maximizing a declared poverty-reduction functional subject to \(K(p,q)\le B\), and estimate whether lower-cost variants preserve durable gains. A factorial design can identify component interactions locally; national-scale effects additionally require variation in coverage and delivery constraints. Define durability using prespecified consumption, assets, and employment measures after support ends. Include next-generation outcomes if the follow-up can identify them. The authors' agenda explicitly asks about essential components and government scalability; this quantitative objective and its constraints are an adapted formulation, not a published optimization theorem.

\textbf{Definitions and admissible scope.} Coverage \(q\) requires a specified recipient-selection rule, not only a fraction. Define the full regime as component doses \(p\), coverage assignment, implementer, start and end dates, and fidelity measures. Let baseline eligible households have weights \(\mu_i\), with exits tracked and household formation assigned an explicit rule. A poverty functional can be \(H_T(p,q)=\sum_i\mu_iP(C_{iT}(p,q)<z)\) for real consumption threshold \(z>0\). Costs include staff, delivery, transfers, and administration under a stated fiscal accounting convention. Maximize \(H_T(0,0)-H_T(p,q)\) over a declared menu subject to budget. At national scale, include nonrecipient welfare and prices; beneficiaries' conditional treatment effects do not determine this population objective. Fix \(T<\infty\), baseline target weights \(\mu_i\), and a real poverty threshold \(z>0\). With \(H_T(p,q)=\sum_i\mu_i\Pr(C_{iT}(p,q)<z)\), the headcount design objective is \(H_T(0,0)-H_T(p,q)\) subject to \(K(p,q)\le B\). The menu includes eligibility and coverage assignment, doses, implementer, and all support dates. If a minimum durable gain is required, add \(E[C_{iT}(p,q)-C_{iT}(0,0)]\ge\underline\tau\) for an explicit \(T\) after each specified support component ends. Identification concerns unknown response laws and feasible delivery constraints; enumeration of known finite-menu poverty values is not itself an open task.
\subsection*{Variants and assumption changes}
\begin{enumerate}
\item Component dose (source-motivated): factorially vary assets, coaching frequency, and digital/group delivery. Test a prespecified noninferiority margin for consumption and poverty at a fixed cost ceiling after support ends.
\item New populations and state delivery (source-motivated): change from rural NGO pilots to urban, refugee, or fragile-state populations and government implementers. Identify selection and delivery effects separately from changes in coverage and market conditions.
\end{enumerate}
\subsection*{References and source audit}
\textbf{Support and access.} Primary institutional agenda checked. Component doses, delivery and population adaptations are explicit questions; formal objectives and noninferiority criteria are editorial. Page is a living resource. \textbf{Locator:} What is Next for Graduation Evidence, paragraphs 2--5.
\begin{enumerate}\raggedright
\item\hypertarget{definitions:source:30007653:1}{} Innovations for Poverty Action. 2025. \emph{The Ultra-Poor Graduation Approach: Research and Learning Agenda for the Next Wave of Graduation Programs}. Section "What’s Next for Graduation Evidence", paragraphs 2–5.
\url{https://poverty-action.org/graduation-approach}.
\end{enumerate}

\subsection*{Common conventions: Source mathematical conventions}

These conventions apply to each entry unless explicitly replaced.

\textbf{Models and identification.} A primitive model \(m\) specifies all probability laws, feasible choices, information, equilibrium and measurement rules used by its target. The admissible class \(\mathcal M\) is fixed independently of outcomes. If \(P_O(m)\) is its observable probability law and \(\tau(m)\) the target, its identified set at observable law \(P\) is
\[
 \mathcal I_\tau(P)=\{\tau(m):m\in\mathcal M,\ P_O(m)=P\}.
\]
Point identification means that this set is a singleton. An empty set signals incompatibility of data and assumptions. Coordinate bounds need not describe the joint set. Bounds are sharp only if every retained target is attainable or arbitrarily approachable by compatible models. Finite bounds require explicit restrictions on unobserved quantities; unrestricted confounding, hidden wealth, extrapolation or arbitrary equilibrium selection can yield uninformative sets.

\textbf{Causal interventions.} A potential outcome \(Y_i(p)\) is the outcome under a complete policy regime \(p\), including timing, financing, information and market spillovers. The target population and units are fixed before treatment. With interference, use the complete assignment vector or a justified exposure mapping; individual no-interference notation alone cannot identify an economy-wide intervention. A difference of mean potential outcomes does not require the cross-world joint distribution. Conditioning on post-treatment migration, survival, enrollment or employment changes the estimand and needs separate justification.

\textbf{Statistical statements.} An estimator, confidence set, forecast or test specifies a sampling law, model class and loss/coverage criterion. Uniform \((1-\alpha)\)-coverage means \(\inf_{m\in\mathcal M}\Pr_m\{\tau(m)\in C_n\}\geq1-\alpha\), or an explicitly asymptotic version. The whole parameter space trivially has coverage, so informativeness requires a stated width, risk or power objective. A forecasting improvement is relative to a named benchmark, information set and evaluation population.

\textbf{Equilibrium and optimization.} An equilibrium comprises feasible plans, optimality, beliefs where needed, budget constraints and all market/resource clearing. The primitives or a stated benchmark define each correspondence \(\mathcal E(p;m)\). Nonemptiness is assumed or proved, never inferred just from compact policy sets. A selected-equilibrium target uses a declared selection rule; a robust target quantifies over all permitted equilibria. Use suprema or infima unless attainment is established. An \(\varepsilon\)-optimal policy is feasible and within \(\varepsilon\) of the finite optimum value. Report an empty feasible set separately.

\textbf{Welfare and accounting.} Normative utility, interpersonal normalization, social weights, numeraire, discounting and the use of public receipts are primitives. Behavioral utility need not coincide with the welfare criterion. Household welfare includes ownership income and rents once; adding profits again to compensating variation generally double-counts them. A monetary equivalent is defined by an explicit transfer and price convention, with monotonicity and range conditions. Average effects, mechanism contributions, transfers and welfare losses are different targets. Infinite sums require convergence and dynamic feasible plans require terminal or no-Ponzi conditions.

\textbf{Source versions.} A working-paper date, revision date and journal publication year are distinguished. A DOI for a journal version is not automatically the DOI of an earlier manuscript. Locators refer to the stated version and printed pages where specified. A metadata-only check does not verify a claim about a full-text conclusion. Every entry's source subsection narrows the provenance claim accordingly.
% END ECONOMICS STATEMENT
\end{document}
